% Generated draft: requires mathematical and source review.
\documentclass[11pt]{article}
\usepackage[margin=25mm]{geometry}
\usepackage{fontspec}
\setmainfont{Latin Modern Roman}
\usepackage{amsmath,amssymb,mathtools,unicode-math}
\setmathfont{Latin Modern Math}
\usepackage{enumitem,needspace,xurl,hyperref}
\setlength{\parindent}{0pt}
\setlength{\parskip}{5pt}
\setlength{\emergencystretch}{3em}
\newcommand{\N}{\mathbb{N}}
\newcommand{\Z}{\mathbb{Z}}
\newcommand{\Q}{\mathbb{Q}}
\newcommand{\R}{\mathbb{R}}
\newcommand{\C}{\mathbb{C}}
\newcommand{\F}{\mathbb{F}}
\newcommand{\A}{\mathbb{A}}
\newcommand{\PP}{\mathbb{P}}
\newcommand{\E}{\mathbb{E}}
\newcommand{\eps}{\varepsilon}
\newcommand{\dd}{\,\mathrm{d}}
\begin{document}
\section*{UnsolvedMath ID: 30001593}
\textbf{Source problem code:} catalogue entry
\subsection{Definitions and mathematical statement}
The supplied record concerns equivariant Schr\"odinger maps on the spatial domain $\R^2$. It does not specify the target manifold, the precise evolution equation, the equivariance action, or the admissible values of the equivariance index $m$. These are unresolved parts of the formulation; in particular, no particular target or range of $m$ is assumed here.

Conditionally on fixing those data, let $\mathcal Q_m$ denote the harmonic-map soliton family called $Q_m$ in the record. For a specified target metric, a harmonic map is a map $Q$ whose tension field vanishes, $\tau(Q)=0$. The record describes $\mathcal Q_m$ as noncompact but does not specify its parameters or the topology in which noncompactness is intended.

If the maps are realized in a fixed common Euclidean ambient space, define the homogeneous Sobolev distance to this family by
\[
 d_m(v):=\inf_{Q\in\mathcal Q_m}\|v-Q\|_{\dot H^1(\R^2)},
 \qquad
 \|v-Q\|_{\dot H^1(\R^2)}^2
 :=\int_{\R^2}|\nabla(v-Q)(x)|^2\,\dd x.
\]
This definition is conditional on an ambient realization and its norm; neither is supplied. If another definition of the map-valued $\dot H^1$ distance is intended, it must be specified instead.

\textbf{Problem.} For equivariant initial data $u_0$ close to $\mathcal Q_m$ in $\dot H^1$, characterize the global dynamics of the corresponding Schr\"odinger map evolution. In particular, determine whether finite-time blow-up can occur while the solution remains close to $\mathcal Q_m$.

More explicitly, after fixing a solution class and a positive closeness tolerance $\varepsilon$, the latter question asks whether there can exist a solution $u$ on its maximal forward interval $[0,T_+)$ such that
\[
 0<T_+<\infty,
 \qquad
 d_m(u(0))<\varepsilon,
 \qquad
 d_m(u(t))<\varepsilon\quad\text{for every }0\leq t<T_+,
\]
with no continuation through $T_+$ in that solution class. The approximating member of $\mathcal Q_m$ may depend on $t$; closeness to the family does not require closeness to one fixed soliton. The record does not specify the permitted tolerances, whether such behavior is sought for arbitrarily small tolerances, or whether initial and subsequent closeness should use different thresholds. Thus the displayed conditions explain the finite-time scenario without imposing an unsupported smallness quantifier. The solution class, blow-up criterion, spatial boundary conditions, and precise alternatives required for a characterization of global dynamics also remain unspecified.
\subsection{Short English statement}
Characterize the evolution of equivariant Schr\"odinger maps on $\R^2$ starting near a noncompact harmonic-map soliton family. In particular, decide whether a solution can develop a finite-time singularity while staying near that family throughout its lifespan. The target, equivariance class, closeness thresholds, and meaning of blow-up must be specified before this becomes a fully determined mathematical problem.
\subsection{Sources}
\begin{itemize}
\item Supplied local UnsolvedMath registry record: permanent ID $30001593$, exact source problem code \texttt{catalogue entry}, title ``Near-Soliton Dynamics of Equivariant Schr\"odinger Maps,'' in the Oberwolfach Reports Open Problems collection. No individual proposer is supplied; the record gives proposed year $2010$.
\item Supplied report citation: \emph{Nonlinear Waves and Dispersive Equations} (2011), Oberwolfach Reports, pp.~2393--2463, DOI \texttt{10.4171/owr/2010/41}. Supplied source and literature link: \url{https://doi.org/10.4171/owr/2010/41}. The background describes these reports as documenting mathematical workshops and their open-problem discussions. No problem-specific page, speaker attribution, or verbatim report passage is supplied.
\item Additional supplied literature link: \url{https://arxiv.org/abs/1009.1608}. The record supplies no title, authors, or precise theorem reference for this item.
\item Literature assessment dated exactly \texttt{2026-08-21}: the supplied curation labels the problem ``partially solved'' and states that near-soliton behavior is known in several equivariant classes, including instability and global solutions near the one-equivariant ground state, while the broad extracted question remains unresolved. This is supplied dated curation, not an independently verified literature assessment; the background explicitly requests independent verification of time-sensitive status and cited claims.
\item The registry labels the statement \texttt{corrected\_verified} and reports a manual semantic, typographic, and LaTeX review of its cleanup. These are supplied editorial labels and do not establish that the missing mathematical specifications have been resolved.
\end{itemize}
\end{document}
